% Part: computability
% Chapter: recursive-functions
% Section: normal-form

\documentclass[../../../include/open-logic-section]{subfiles}

\begin{document}

\olfileid{cmp}{rec}{nft}
\olsection{د نورمال شکل قضيه}

\begin{thm}[د کليني د نورمال شکل قضيه]
\ollabel{thm:kleene-nf}
يوه بنسټيزه بازګشتي اړيکه $T(e, x, s)$ او يوه بنسټيزه بازګشتي تابع
$U(s)$ شته چې دا خاصيت لري: که $f$ هره جزوي بازګشتي تابع وي، نو د
کوم~$e$ دپاره
\[
f(x) \simeq U(\umin{s}{T(e, x, s)})
\]
د هر~$x$ دپاره پوره کېږي۔
\end{thm}

\begin{explain}
د نورمال شکل د قضيې ثبوت اوږد دے، خو بنسټيزه مفکوره ئې ساده ده۔ هره
جزوي بازګشتي تابع يو \emph{شاخص}~$e$ لري؛ په شهودي ډول دا هغه عدد دے
چې د تابع پروګرام يا تعريف کوډوي۔ که $f(x) \fdefined$ وي، محاسبه په
منظم ډول ثبتېدے او په کوم عدد~$s$ کوډېدے شي، او دا حقيقت، چې $s$ پر
ننوت~$x$ د~$f$ محاسبه کوډوي، يوازې د~$x$ او تعريف~$e$ په کارولو سره
په بنسټيز بازګشتي ډول کتل کېدے شي۔ نو اړيکه~$T$، يعنې ``د شاخص~$e$
تابع د ننوت~$x$ دپاره يوه محاسبه لري، او $s$ دا محاسبه کوډوي''،
بنسټيزه بازګشتي ده۔ د محاسبې د بشپړ ريکارډ~$s$ په ورکړې سره د~$s$
``پايله'' د~$f(x)$ قيمت دے، او هغه هم له~$s$ څخه په بنسټيز بازګشتي
ډول تر لاسه کېدے شي۔

د نورمال شکل قضيه ښيي چې د هرې جزوي بازګشتي تابع د تعريف دپاره يوازې
يوه بې‌حده پلټنه پکار ده۔ په بنسټيز ډول د ټولو عددونو له منځه پلټنه
کوو تر څو داسې عدد ومومو چې د شاخص~$e$ د تابع محاسبه پر ننوت~$x$
کوډوي۔ عددونه~$e$ د جزوي بازګشتي تابعو د ``نومونو'' په توګه کارولے
شو، او د قضيې په معادله تعريف شوې تابع~$f$ ته $\cfind{e}$ ليکلے شو۔
پام وکړئ چې يوه جزوي بازګشتي تابع تر يوه ډېر شاخص لرلے شي---په حقيقت
کښې هره جزوي بازګشتي تابع بې‌شمېره ډېر شاخصونه لري۔
\end{explain}
\end{document}
